\documentclass[12pt,a4paper]{article}

% ===== Technical Report format / Chinese LaTeX setting =====
% 建議使用 XeLaTeX 編譯。
\usepackage[margin=1in]{geometry}
\usepackage{fontspec}
\usepackage{xeCJK}

\setmainfont{Times New Roman}

\setCJKmainfont{AR PL UKai TW}
\setCJKsansfont{AR PL UKai TW}
\setCJKmonofont{AR PL UKai TW}

\usepackage{setspace}
\usepackage{cite}
\usepackage{amsmath,amssymb,amsfonts}
\usepackage{booktabs}
\usepackage{array}
\usepackage{tabularx}
\usepackage{longtable}
\usepackage{graphicx}
\usepackage{tikz}
\usepackage{xcolor}
\usepackage{url}
\usepackage{tocloft}
\usepackage[hidelinks]{hyperref}
\usepackage{fancyhdr}
\usepackage{titlesec}
\usepackage{caption}
\usepackage{float}
\usepackage{tocloft}
\usepackage{hyperref}
\usepackage{float}
\usepackage{placeins}

% ===== Table of contents dot leaders =====
\renewcommand{\cftsecleader}{\cftdotfill{\cftdotsep}}
\renewcommand{\cftsubsecleader}{\cftdotfill{\cftdotsep}}
\renewcommand{\cftsubsubsecleader}{\cftdotfill{\cftdotsep}}

% ===== Make TOC entries normal weight =====
\renewcommand{\cftsecfont}{\normalfont}
\renewcommand{\cftsecpagefont}{\normalfont}
\renewcommand{\cftsubsecfont}{\normalfont}
\renewcommand{\cftsubsecpagefont}{\normalfont}

\usetikzlibrary{arrows.meta,positioning,fit,calc}

\newcolumntype{Y}{>{\raggedright\arraybackslash}X}
\newcolumntype{C}{>{\centering\arraybackslash}p{0.15\textwidth}}
\newcommand{\etal}{\textit{et al.}}
\newcommand{\ie}{i.e., }
\newcommand{\eg}{e.g., }

\renewcommand{\figurename}{圖}
\renewcommand{\tablename}{表}
\renewcommand{\refname}{參考文獻}
\renewcommand{\contentsname}{目錄}
\renewcommand{\listfigurename}{圖目錄}
\renewcommand{\listtablename}{表目錄}


\titleformat{\section}{\Large\bfseries}{\thesection}{1em}{}
\titleformat{\subsection}{\large\bfseries}{\thesubsection}{1em}{}
\titleformat{\subsubsection}{\normalsize\bfseries}{\thesubsubsection}{1em}
{}

\pagestyle{fancy}
\fancyhf{}
\lhead{Data Center Thermal Prediction Technical Report}
\rhead{TR-2026-0617}
\cfoot{\thepage}
\renewcommand{\headrulewidth}{0.4pt}
\renewcommand{\footrulewidth}{0.4pt}

% 如果圖片尚未上傳到 Overleaf，先產生占位框，避免編譯失敗。
\newcommand{\safeincludegraphics}[2][]{%
  \IfFileExists{#2}{\includegraphics[#1]{#2}}{%
    \fbox{\begin{minipage}{0.88\textwidth}\centering\vspace{1cm}
    圖片檔案尚未上傳：\\\texttt{#2}\\
    \vspace{1cm}\end{minipage}}%
  }%
}

\setlength{\parindent}{2em}
\setlength{\parskip}{0.4em}
\onehalfspacing

% ===== Table of contents dot leaders =====
\renewcommand{\cftsecleader}{\cftdotfill{\cftdotsep}}
\renewcommand{\cftsubsecleader}{\cftdotfill{\cftdotsep}}
\renewcommand{\cftsubsubsecleader}{\cftdotfill{\cftdotsep}}

% ===== Section title normal weight =====
\titleformat{\section}
  {\normalfont\Large}
  {\thesection}{1em}{}

\titleformat{\subsection}
  {\normalfont\large}
  {\thesubsection}{1em}{}

\titleformat{\subsubsection}
  {\normalfont\normalsize}
  {\thesubsubsection}{1em}{}

\begin{document}
\setstretch{1.5}
% ===== Title Page =====
\begin{titlepage}
\thispagestyle{empty}
\centering
\vspace*{1cm}

\vspace*{\fill}

{\Huge 資料中心熱監測、機櫃層級溫度預測\\
與功率感知利用之 AI/ML 模型技術報告\par}
\vspace{1.5cm}

\begin{tabular}{ll}
\toprule
\textbf{Report ID} & TR-2026-0617 \\
\textbf{Version} & v1.0 \\
\textbf{Release Date} & 2026.06.17 \\
\textbf{Report Type} & Technical Survey Report \\
\bottomrule
\end{tabular}

\vspace*{\fill}

\vspace{1.5cm}


\vfill
\end{titlepage}

\pagenumbering{roman}

% ===== Executive Summary =====
\section*{Executive Summary}
\addcontentsline{toc}{section}{Executive Summary}
現代雲端資料中心的溫度預測相當困難，因為熱行為會受到工作負載、功率、冷卻、氣流、機櫃配置、設備異質性以及感測器位置等多個因素共同影響。因此，當同一個機櫃可能包含不同類型的伺服器、不同的 U 位置分布，以及不完整的感測覆蓋範圍時，傳統的查找表方法便不足以應付。本調查報告整理現有的人工智慧與機器學習模型，這些模型能夠學習多因素之間的熱相關性，並評估它們如何支援溫度監測、CPU/GPU 過熱預防，以及功率消耗監測，以提升機櫃利用率。已回顧的文獻顯示，XGBoost 與 LightGBM 是以結構化遙測資料進行伺服器與機櫃溫度預測的強力候選模型；循環神經網路則能支援具時間相依性的熱行為；U-Net 與隱式神經表示法可支援由稀疏觀測到完整熱場的重建；數位孿生能提供安全的視覺化與策略評估環境；而強化學習則可支援工作負載與冷卻控制。根據已驗證的 CortexDC 遙測參考資料，本報告提出一個保守的機櫃層級監測架構，使用 Redfish 伺服器遙測資料、PDU 量測資料、環境感測器與機櫃中繼資料作為已證實可用的資料來源。所提出的架構將伺服器層級預測與機櫃層級預測分開，並明確加入空間特徵，例如機櫃 ID、U 位置、鄰近機櫃/伺服器溫度，以及機櫃功率分布。

\section*{Keywords}
\addcontentsline{toc}{section}{Keywords}
資料中心熱管理，機櫃層級溫度預測，XGBoost，LightGBM，Redfish，PDU 監測，感測器配置，熱重建，功率監測

\newpage
\clearpage
\tableofcontents

\clearpage
\phantomsection
\addcontentsline{toc}{section}{圖目錄}
\listoffigures

\clearpage
\phantomsection
\addcontentsline{toc}{section}{表目錄}
\listoftables

\clearpage
\newpage

% ===== List of Abbreviations =====
\clearpage
\phantomsection
\addcontentsline{toc}{section}{List of Abbreviations}
\section*{List of Abbreviations}

表~\ref{tab:abbreviations} 列出本技術報告中使用的主要縮寫與其定義。
\begin{longtable}{p{0.22\textwidth} p{0.68\textwidth}}
\caption{Definitions of Abbreviations（縮寫表定義）}
\label{tab:abbreviations}\\

\toprule
Abbreviation & Definition \\
\midrule
\endfirsthead

\toprule
Abbreviation & Definition \\
\midrule
\endhead

AI/ML & Artificial Intelligence / Machine Learning，人工智慧與機器學習 \\

BMC & Baseboard Management Controller，主機板管理控制器 \\

CPU & Central Processing Unit，中央處理器 \\

GPU & Graphics Processing Unit，圖形處理器 \\

PDU & Power Distribution Unit，配電單元 \\

Redfish & 伺服器管理與遙測資料存取介面 \\

XGBoost & Extreme Gradient Boosting，梯度提升決策樹模型 \\

LightGBM & Light Gradient Boosting Machine，輕量化梯度提升模型 \\

GRU & Gated Recurrent Unit，門控循環單元 \\

LSTM & Long Short-Term Memory，長短期記憶網路 \\

U-Net & 具有編碼器與解碼器結構的卷積神經網路 \\

INR & Implicit Neural Representation，隱式神經表示法 \\

DQN & Deep Q-Network，深度 Q 網路 \\

RL & Reinforcement Learning，強化學習 \\

PUE & Power Usage Effectiveness，電源使用效率 \\

CFD & Computational Fluid Dynamics，計算流體力學 \\

CRAC & Computer Room Air Conditioner，電腦機房空調 \\

\bottomrule
\end{longtable}



\newpage
\pagenumbering{arabic}

\section{緒論}

雲端資料中心會將大部分消耗的電能轉換成熱能。如果熱行為無法被準確監測，局部熱點可能會降低硬體可靠度、觸發 CPU/GPU 降頻、增加風扇能耗，並迫使資料中心使用過度的機房層級冷卻。因此，熱感知管理同時需要準確的溫度預測與可執行的功率監測。

本報告的目標問題，是由多樣化的機櫃設備配置與不同感測器位置所造成的溫度預測困難。在真實機櫃中，伺服器型號、U 位置、有效功率、風扇行為、進氣溫度，以及鄰近設備可能會隨時間而不同。這些影響是非線性且彼此耦合的，因此靜態查找表無法很好地泛化到新的機櫃配置或工作負載狀態。近期的回顧研究也強調，資料中心熱模型必須同時考慮工作負載、冷卻、氣流與熱分布，而不能將它們視為彼此獨立的因素~\cite{lin_review_2024}。

本調查的目標，是找出能夠處理這類多因素相關性的 AI/ML 模型，並將文獻轉化為適用於 CortexDC 監測資料中心環境的實務架構。本報告聚焦於三個操作成果：溫度監測、CPU/GPU 過熱預防，以及用於機櫃利用率的功率消耗監測。本報告並不宣稱有新的實驗結果，而是根據 CortexDC 參考文件中記載的遙測資料，提出一個保守的架構。

本報告的主要貢獻可整理如下：
\begin{itemize}
\item 提出一個系統架構，說明伺服器、PDU、機櫃/環境，以及 CortexDC/iBox 資料如何被收集與處理。
\item 整理 AI/ML 模型家族與伺服器層級預測、機櫃層級預測、熱重建、感測器配置與最佳化之間的對應關係。
\item 提出一個機櫃層級架構，將伺服器預測與機櫃預測分開，並加入機櫃層級建模所需的空間特徵。
\item 建立一個保守的遙測可用性表格，只使用 Redfish、PDU 與 CortexDC 機櫃文件中已證實存在的欄位。
\end{itemize}

\section{系統架構與資料來源}

圖~\ref{fig:system_architecture} 顯示在套用詳細功能架構之前所考慮的環境。實體層包含伺服器、機櫃、配電單元（PDU）與環境感測器。伺服器遙測資料會透過 Redfish 從主機板管理控制器（BMC）中收集。PDU 資料則依照裝置能力，透過 SNMP、Modbus TCP 或 Redfish 收集。CortexDC/iBox 會儲存近期與歷史遙測資料，並將資料提供給特徵工程、第一層 XGBoost/LightGBM 預測、後續的數位孿生熱圖映射、未來的強化學習決策最佳化，以及預防機制使用。

\begin{figure}[htbp]
\centering
\safeincludegraphics[width=\textwidth]{figures/system_architecture_cortexdc_v2.png}
\caption[所提出架構之系統架構]{所提出架構之系統架構}
\label{fig:system_architecture}
\end{figure}


\subsection{CortexDC 中已驗證的遙測資料}

遙測層必須採取保守方式，因為不是所有理想特徵都能從機櫃文件中取得。CortexDC Redfish 參考文件記載了伺服器功率、溫度感測器、風扇速度、由供應商/OEM 介面暴露的 GPU 指標，以及健康狀態與電源狀態欄位~\cite{cortexdc_redfish_ref}。PDU 參考文件記載了進線/出線電氣量測、總有效功率、有效電能、功率因數、環境感測器與歷史趨勢~\cite{cortexdc_pdu_ref}。CortexDC 機櫃清單文件則記載了機櫃中繼資料、U 位置、Redfish 間隔、伺服器功率、進氣溫度、CPU 溫度，以及若干可能對部分資產不可用的欄位~\cite{cortexdc_datasheet_2026}。

表~\ref{tab:telemetry_availability} 整理本報告宣稱可用的資料。CPU 使用率與其他 Node Exporter 系統指標並未被刻意作為必要特徵，因為機櫃文件經常將它們標記為 N/A。若它們存在於實際部署的遙測串流中，未來可能會有用；但在第一版架構中不應被預設為可用。

\begin{table}[t]
\scriptsize
\centering
\caption{所提出架構採用的保守遙測可用性。}
\label{tab:telemetry_availability}
\footnotesize
\begin{tabularx}{\textwidth}{p{0.17\textwidth} p{0.23\textwidth} p{0.18\textwidth} Y}
\toprule
\textbf{資料群組} & \textbf{已驗證欄位} & \textbf{來源} & \textbf{在所提出架構中的用途} \\
\midrule
伺服器功率 & 目前功率、平均功率、最大功率、最小功率、容量、已分配功率、可用功率、請求功率，或裝置可提供的限制功率 & 伺服器 BMC Redfish~\cite{cortexdc_redfish_ref}、機櫃文件~\cite{cortexdc_datasheet_2026} & 作為伺服器熱源代理、總機櫃功率聚合、功率趨勢與機櫃利用率估計 \\
\addlinespace
伺服器熱資料 & 溫度感測器、進氣溫度、CPU 平均溫度、溫度閾值、健康狀態 & 伺服器 BMC Redfish~\cite{cortexdc_redfish_ref}、機櫃文件~\cite{cortexdc_datasheet_2026} & 伺服器層級溫度預測、機櫃最高溫度與熱點警告 \\
\addlinespace
冷卻/風扇 & 由 BMC 或外部感測器提供的風扇 RPM 與風扇備援/狀態 & 伺服器 BMC Redfish~\cite{cortexdc_redfish_ref}、PDU/環境感測器支援~\cite{cortexdc_pdu_ref} & 冷卻反應特徵、風扇異常訊號與過熱預防建議 \\
\addlinespace
GPU/加速器 & 只有當標準/OEM Redfish 有提供時，才包含 GPU 型號、記憶體、使用率、時脈、功率與溫度 & 伺服器 BMC Redfish~\cite{cortexdc_redfish_ref} & GPU 過熱監測，以及在可用時評估 GPU 熱源貢獻 \\
\addlinespace
PDU 電氣資料 & 進線/出線電壓、電流、有效功率、視在功率、功率因數、有效電能、插座狀態與總有效功率 & PDU SNMP/Modbus
/Redfish ~\cite{cortexdc_pdu_ref} & 機櫃功率驗證、插座層級功率、能源趨勢、過載與低利用率檢查 \\
\addlinespace
環境資料 & 若有安裝感測器，則包含溫度、濕度、氣流、壓力、門磁、煙霧/水浸/振動 & PDU/環境感測器支援~\cite{cortexdc_pdu_ref} & 外部熱環境、冷卻分布與感測器配置評估 \\
\addlinespace
機櫃中繼資料 & 機櫃 ID、資產/伺服器名稱、U 位置、U 大小、伺服器放置順序、線上/離線狀態 & 機櫃清單中繼資料~\cite{cortexdc_datasheet_2026} & 空間特徵工程、鄰近伺服器/機櫃特徵、上/中/下層機櫃功率 \\
\addlinespace
不預設可用 & CPU 使用率、記憶體使用率、網路流量、磁碟 I/O、TCP 連線數 & 機櫃文件顯示許多系統指標為 N/A~\cite{cortexdc_datasheet_2026} & 若未來持續可用，可作為選用特徵；不作為所提出基準模型的必要條件 \\
\bottomrule
\end{tabularx}
\end{table}

\section{熱與功率預測的 AI/ML 模型調查}

表~\ref{tab:literature_mapping} 將已回顧的論文對應到它們所支援的架構層。此表的目的不是對論文排名，而是呈現每一篇來源如何對完整的監測與建議系統有所貢獻。

\begin{table}[t]
\scriptsize
\centering
\caption{已回顧論文與所提出架構之對應。}
\label{tab:literature_mapping}
\footnotesize
\begin{tabularx}{\textwidth}{p{0.23\textwidth} p{0.22\textwidth} Y}
\toprule
\textbf{論文} & \textbf{支援的架構層} & \textbf{在本報告中的用途} \\
\midrule
J. Lin \etal{}~\cite{lin_review_2024} & 文獻基礎與開放議題 & 建立熱模型、AI/ML 預測、灰箱模型、熱感知節能，以及 IT 與冷卻聯合最佳化的背景。 \\
\addlinespace
K. Zhang \etal{}~\cite{zhang_2018} & 伺服器/元件預測 & 支援跨系統元件的執行期溫度預測與熱感知任務配置。 \\
\addlinespace
S. Wang \etal{}~\cite{wang_self_healing_2024} & 時序預測與排程 & 支援以 GRU 為基礎的溫度預測，包含延遲熱行為、鄰近伺服器資訊與自我修復式再訓練。 \\
\addlinespace
I. Nisce \etal{}~\cite{nisce_2023} 與 S. Ilager \etal{}~\cite{ilager_2021} & XGBoost 伺服器/主機預測 & 支援 XGBoost 用於結構化伺服器或主機遙測資料，以及熱感知資源管理。 \\
\addlinespace
J. Lin \etal{}~\cite{lin_ate_2022} & 機櫃/機房預測 & 支援 XGBoost 與 LightGBM 用於氣冷式資料中心熱預測，並使用機櫃功率與冷卻參數。 \\
\addlinespace
K. Sano \etal{}~\cite{sano_2018} & 空間/設施熱建模 & 說明進氣溫度為何會受到氣流、機櫃細節、IT 負載與設施側冷卻影響。 \\
\addlinespace
X. Wang \etal{}~\cite{xwang_sensor_2013} 與 T. R. Andersson \etal{}~\cite{andersson_2023} & 感測器配置 & 支援智慧型或不確定性感知的感測器配置，以改善熱點或未觀測位置的偵測。 \\
\addlinespace
H. Yuan \etal{}~\cite{yuan_2025}、S. Sabathiel \etal{}~\cite{sabathiel_2024}、X. Peng \etal{}~\cite{peng_2022} 與 X. Luo \etal{}~\cite{luo_2024} & 熱重建與數位孿生 & 支援使用 U-Net 類型、物理資訊、數位孿生或隱式神經表示法，從稀疏觀測重建完整熱場。Yuan \etal{} 也展示數位孿生如何連接監測、視覺化、重建與以 DQN 為基礎的冷卻決策。 \\
\addlinespace
M. Zapater \etal{}~\cite{zapater_2015}、C. S. Chan \etal{}~\cite{chan_2019} 與 Y. Ran \etal{}~\cite{ran_2019} & 最佳化與決策支援 & 支援將預測連接到風扇控制、效能感知冷卻、工作負載排程，以及工作負載與冷卻的聯合最佳化。 \\
\bottomrule
\end{tabularx}
\end{table}

\subsection{結構化回歸模型}

XGBoost 與 LightGBM 適合用於結構化遙測資料，因為它們能夠建模非線性關係、容忍混合尺度的特徵，並且能良好處理經過設計的表格型特徵。Nisce \etal{} 使用 XGBoost，根據使用率、溫度與能量相關資料進行 CPU/伺服器層級熱預測，並在雲端伺服器溫度預測中回報了低預測誤差~\cite{nisce_2023}。Ilager \etal{} 比較了多種用於雲端主機溫度預測的機器學習模型，並在熱感知排程系統中使用 XGBoost 作為最強的預測器~\cite{ilager_2021}。Lin \etal{} 評估了用於氣冷式資料中心熱預測的資料驅動模型，並發現 XGBoost 與 LightGBM 在穩態機櫃進氣溫度預測上表現良好~\cite{lin_ate_2022}。

這些研究支持將 XGBoost/LightGBM 作為第一層預測模型，特別是當輸入為結構化遙測資料時，例如功率、進氣溫度、CPU 溫度、風扇速度、機櫃功率與機櫃中繼資料。然而，單靠 XGBoost 並不足以解決完整的機櫃層級問題。它可以預測數值，但無法自動重建空間熱場、推論缺少感測器的區域，或最佳化工作負載/冷卻行動。

\subsection{時序模型}

熱行為具有延遲性：工作負載與功率可能快速變化，但由於熱傳與冷卻慣性的影響，溫度變化較慢。Zhang \etal{} 顯示機器學習可用於預測元件溫度，以支援執行期熱管理，並引導熱感知任務配置~\cite{zhang_2018}。Wang \etal{} 提出一個兩階段自我修復模型，先預測 CPU 使用率，再利用 CPU、風扇、進氣/排氣溫度、CRAC 溫度與鄰近伺服器溫度資料來預測未來伺服器進氣溫度~\cite{wang_self_healing_2024}。ThermoSim 也展示循環神經網路如何支援雲端環境中的熱感知模擬與資源管理~\cite{gill_thermosim_2020}。

對於所提出的架構而言，當歷史遙測資料足夠穩定且密集時，時序模型最有用。如果只有 5 分鐘的 Redfish 收集間隔，簡單的延遲特徵與趨勢特徵可能已足以建立基準模型。若未來能收集更長期的歷史資料，並取得更可靠的工作負載指標，則可考慮 GRU/LSTM 模型。

\subsection{空間建模與感測器配置}

機櫃層級預測需要空間感知。伺服器 U 位置、機櫃位置、垂直機櫃區域、伺服器分布與氣流路徑，都可能改變同一個功率值所造成的熱效應。Sano \etal{} 顯示，詳細的 CFD 建模可以提升資料中心溫度分布預測的準確度，因為氣流與機櫃縫隙會影響進氣溫度~\cite{sano_2018}。X. Wang \etal{} 指出，感測器配置會影響熱伺服器偵測，而由 CFD 引導的配置方式比均勻配置更能偵測過熱~\cite{xwang_sensor_2013}。Andersson \etal{} 研究了具有不確定性感知的環境感測器配置，並使用卷積式高斯神經過程；雖然其應用場景並非資料中心，但其概念可用於判斷新增量測位置如何降低不確定性~\cite{andersson_2023}。

這些論文支持機櫃層級預測必須加入空間特徵，而不能只使用伺服器遙測資料。主要空間特徵包括機櫃 ID、U 位置、U 大小、伺服器放置順序、垂直機櫃區域、加權 U 位置，以及上/中/下層機櫃功率。若鄰近熱資料可用，也可以加入鄰近伺服器或鄰近機櫃溫度作為選用的空間背景特徵。

\subsection{熱場重建}

當感測器稀疏時，單一預測數值不足以呈現機櫃或機房的熱狀態。熱場重建方法可以從有限觀測推論完整熱圖。Yuan \etal{} 提出一個數位孿生原型，使用感測器資料、伺服器功率與 CRAC 溫度作為輸入；U-Net 代理模型用於重建 2D 溫度場，而 DQN 決策代理則根據 PUE 降低、溫度均勻性與熱點抑制來建議冷卻決策~\cite{yuan_2025}。這使數位孿生不只是視覺化工具，也成為連接預測與安全最佳化的重要橋梁。Sabathiel \etal{} 使用神經網路架構，從稀疏量測重建未知材料組成系統中的穩態溫度場~\cite{sabathiel_2024}。Peng \etal{} 使用基於分區建模的自適應 U-Net 與 MLP，從有限觀測點重建溫度場~\cite{peng_2022}。Luo \etal{} 則使用隱式神經表示法，從稀疏觀測重建連續場~\cite{luo_2024}。

對於目標機櫃使用案例而言，熱重建應該是基準遙測預測被驗證後的第二階段延伸。U-Net 或隱式神經表示法需要空間網格、感測器座標，以及能近似機櫃或機房熱場的訓練資料。

\subsection{最佳化與控制}

預測只有在連接到監測或控制決策時，才具有操作價值。Zapater \etal{} 顯示，伺服器風扇控制應考慮漏電功率與風扇功率之間的取捨~\cite{zapater_2015}。Chan \etal{} 加入效能面向，指出冷卻風扇決策可能影響硬碟效能與應用程式延遲~\cite{chan_2019}。Ran \etal{} 提出 DeepEE，這是一個強化學習架構，用於聯合最佳化資料中心能源效率中的工作排程與冷卻控制~\cite{ran_2019}。

這些研究建議採用分階段控制路徑。XGBoost/LightGBM 應被提出作為第一個實務預測層，因為它們可以在不改變真實系統的情況下，使用歷史結構化遙測資料訓練。強化學習則應在安全模擬器、數位孿生或具規則限制的環境可用後，作為未來的決策最佳化延伸。在此分階段設計中，數位孿生可以重播或模擬熱反應、評估候選行動，並降低在真實機櫃上學習不安全控制策略的風險。直接致動行為，例如工作負載遷移、風扇控制或冷卻設定點控制，應先定義控制介面與安全限制。

\newpage
\section{所提出的機櫃層級預測與預防架構}

圖~\ref{fig:functional_framework} 顯示所提出的功能方塊圖。它修正了早期架構，從實體資料來源開始，並明確呈現每個元件的輸入、處理步驟與輸出。此架構包含七個階段：資料來源、遙測資料擷取、資料可用性過濾、特徵工程、伺服器層級與機櫃層級預測、熱狀態映射/數位孿生重建，以及以預防為導向的決策支援。

\begin{figure}[H]
\centering
\safeincludegraphics[width=\textwidth]{figures/functional_framework_imagegen.png}
\caption[修正版功能方塊圖]{修正版功能方塊圖}
\label{fig:functional_framework}
\end{figure}


\subsection{輸入與輸出定義}

假設機櫃 $r$ 在時間 $t$ 包含作用中的伺服器 $s \in S_r$。基準機櫃特徵向量定義為
\begin{equation}
\mathbf{x}_{r,t} =
[ \mathbf{p}_{r,t}, \mathbf{q}_{r,t}, \mathbf{h}_{r,t}, \mathbf{z}_{r,t}, \mathbf{d}_{r,t} ],
\end{equation}
其中 $\mathbf{p}_{r,t}$ 包含功率特徵，$\mathbf{q}_{r,t}$ 包含熱特徵，$\mathbf{h}_{r,t}$ 包含健康狀態/風扇特徵，$\mathbf{z}_{r,t}$ 包含空間機櫃特徵，而 $\mathbf{d}_{r,t}$ 包含時間差分或延遲數值。

對於機櫃 $r$ 中的伺服器 $s$，伺服器層級輸入向量可寫為
\begin{equation}
\mathbf{x}_{s,t}^{srv} =
[P_{s,t}, T^{cpu}_{s,t}, T^{gpu}_{s,t}, T^{in}_{s,t}, F_{s,t}, H_{s,t}],
\end{equation}
其中 $P_{s,t}$ 是伺服器功率，$T^{cpu}_{s,t}$ 是 CPU 溫度，$T^{gpu}_{s,t}$ 是 GPU 溫度（若可用），$T^{in}_{s,t}$ 是進氣溫度，$F_{s,t}$ 是風扇速度（若有暴露），而 $H_{s,t}$ 是健康狀態或電源狀態指標。

基準機櫃層級預測問題為
\begin{equation}
\hat{\mathbf{y}}_{r,t+\Delta} = f(\mathbf{x}_{r,t}),
\end{equation}
其中 $\hat{\mathbf{y}}_{r,t+\Delta}$ 可包含預測的機櫃平均溫度、機櫃最高溫度、下一個時間區間的熱點風險，或機櫃利用狀態。時間範圍 $\Delta$ 應與可用的資料收集間隔與操作反應時間相匹配。

\subsection{特徵工程}

基準特徵群組如下：
\begin{itemize}
\item \textbf{功率特徵：} 總機櫃功率、平均伺服器功率、最大伺服器功率、PDU 總有效功率、插座有效功率、有效電能，以及功率因數。
\item \textbf{熱特徵：} 平均 CPU 溫度、最大 CPU 溫度、平均進氣溫度、最大進氣溫度，以及外部環境溫度。
\item \textbf{冷卻/健康特徵：} 若可取得，包含風扇速度、風扇狀態、伺服器健康狀態、電源狀態，以及 GPU 溫度/功率。
\item \textbf{空間特徵：} 機櫃 ID、U 位置、U 大小、上/中/下機櫃區域、加權 U 位置、作用中伺服器數量、鄰近伺服器溫度，以及在可用時的鄰近機櫃溫度。
\item \textbf{時間特徵：} 溫度變化率、功率變化率、移動平均、移動最大值，以及延遲的溫度/功率值。
\end{itemize}

機櫃區域特徵可直接處理建立機櫃層級空間效應模型的需求。主要空間特徵包含機櫃 ID、U 位置、U 大小、伺服器放置順序、垂直機櫃區域、加權 U 位置，以及上/中/下層機櫃功率。當鄰近熱資料可用時，可以加入鄰近伺服器或鄰近機櫃溫度，作為選用的空間背景特徵。例如，上/中/下層機櫃功率可以顯示熱源是垂直集中，還是均勻分布。

\begin{table}[H]
\scriptsize
\centering
\caption{機櫃層級熱預測的候選特徵。}
\label{tab:candidate_features}
\footnotesize
\begin{tabularx}{\textwidth}{p{0.18\textwidth} p{0.29\textwidth} p{0.21\textwidth} Y}
\toprule
\textbf{特徵群組} & \textbf{範例特徵} & \textbf{主要來源} & \textbf{目的} \\
\midrule
伺服器功率 & 目前功率、平均功率、最大功率、最小功率、PSU 數量、可用時的 GPU 功率 & Redfish / BMC & 估計伺服器產熱量與其對機櫃功率的貢獻。 \\
\addlinespace
伺服器熱資料 & CPU 溫度、可用時的 GPU 溫度、進氣溫度、溫度感測器狀態、風扇速度 & Redfish / BMC & 描述伺服器熱狀態與冷卻反應。 \\
\addlinespace
PDU 功率 & 總有效功率、插座有效功率、電壓、電流、功率因數、有效電能、插座狀態 & PDU 經由 SNMP、Modbus 或 Redfish & 提供機櫃層級電氣負載與功率驗證。 \\
\addlinespace
機櫃中繼資料 & 機櫃 ID、U 位置、U 大小、伺服器順序、垂直區域、線上/離線狀態 & 清單 / 機櫃中繼資料 & 為機櫃層級預測加入空間背景。 \\
\addlinespace
環境資料 & 外部溫度、濕度、氣流、壓力、若有安裝則包含門磁 & PDU/環境感測器 & 加入冷卻與機房環境資訊。 \\
\addlinespace
時間趨勢 & 移動平均、移動最大值、延遲溫度、延遲功率、溫度變化率、功率變化率 & 歷史遙測資料 & 捕捉延遲熱反應與短期風險。 \\
\addlinespace
選用工作負載 & CPU 使用率、記憶體使用率、網路流量、磁碟 I/O、TCP 連線數 & OS 遙測 / Node Exporter，若可用 & 描述工作負載強度，但除非已從遙測紀錄驗證可用，否則不應列為必要特徵。 \\
\bottomrule
\end{tabularx}
\end{table}

\FloatBarrier

\subsection{預測目標}

此架構將伺服器層級與機櫃層級預測目標分開：
\begin{itemize}
\item \textbf{伺服器層級目標：} 下一時刻 CPU 溫度、下一時刻進氣溫度、當 GPU 遙測存在時的 GPU 過熱風險，以及伺服器熱點風險。
\item \textbf{機櫃層級目標：} 機櫃平均溫度、機櫃最高溫度、下一區間機櫃溫度、熱點風險分數，以及機櫃功率利用狀態。
\end{itemize}

XGBoost/LightGBM 應作為這兩個層級的第一候選模型，因為已驗證的遙測資料為表格型資料，並包含經設計的空間特徵。若收集到足夠的歷史軌跡，之後可再評估 GRU/LSTM 模型。U-Net/INR 重建與數位孿生視覺化應在感測器座標與熱圖訓練資料準備完成後再進行評估。強化學習則應在數位孿生或具規則限制的模擬器能安全測試工作負載與冷卻行動之後再評估。

\section{預防與決策機制}

\subsection{用於預防的熱狀態映射}

由於監測功能已由遙測平台處理，因此所提出的貢獻應聚焦於將監測資料轉換成可預防的熱狀態與行動候選項。第一層預防機制應將原始遙測資料映射為四種工程視角：
\begin{itemize}
\item \textbf{機櫃區域圖：} 將伺服器功率、CPU 溫度、進氣溫度、風扇速度與伺服器狀態聚合到上、中、下機櫃區域。
\item \textbf{熱風險圖：} 將預測的伺服器與機櫃溫度轉換為風險等級，例如正常、警告、危急與不穩定趨勢。
\item \textbf{工作負載狀態圖：} 將相似的工作負載或功率--溫度模式分組為操作狀態。當工作負載軌跡可用時，可使用分群方法進行此映射，類似效能感知冷卻控制中使用的工作負載狀態建模~\cite{chan_2019}。
\item \textbf{數位孿生狀態圖：} 當累積足夠空間資料後，可將稀疏遙測資料與預測溫度投影到重建的機櫃或機房熱場中。
\end{itemize}

此映射層將系統目標從「觀察溫度」轉換為「辨識可預防的熱風險」。例如，某個機櫃的絕對溫度可能仍在安全範圍內，但如果其溫度斜率上升、上層區域功率集中，或預測會出現熱點，仍然需要被標記為風險。這些映射後的狀態會成為規則式預防、數位孿生評估與未來強化學習策略的輸入。

\subsection{CPU/GPU 過熱預防機制}

過熱預防應採用主動方式，而不只是依賴閾值。閾值警報通常是在溫度已經偏高後才偵測到問題，而預測模型可以在下一個時間區間變得危險之前發出警告。基準預防機制應包含：
\begin{itemize}
\item 使用 XGBoost 或 LightGBM 估計下一區間的 CPU/進氣溫度；
\item 將預測溫度與風扇速度、目前功率趨勢結合；
\item 僅在 Redfish 有暴露 GPU 溫度與 GPU 功率時，才將它們視為選用特徵；
\item 當預測風險仍偏高時，建議工作負載降低、工作負載遷移、降頻，或硬體檢查；
\item 在控制介面與安全限制尚未被定義之前，避免直接控制風扇或工作負載。
\end{itemize}

\subsection{功率感知的機櫃利用率與過載預防}

功率遙測資料不應只用於顯示耗電量，也應用於預防低效率或不安全的機櫃利用。PDU 有效功率提供機櫃層級的電氣狀態，而伺服器 Redfish 功率則可指出哪些資產對熱生成貢獻最大。預防機制應計算：
\begin{itemize}
\item \textbf{熱風險調整後的利用率：} 當預測最高溫度或熱點風險增加時，機櫃利用率應降低。
\item \textbf{區域功率不平衡：} 應檢查上/中/下層機櫃功率，以避免熱源過度集中。
\item \textbf{伺服器貢獻排名：} 高功率且高預測溫度的伺服器，應優先被考慮進行工作負載降低或遷移。
\item \textbf{低利用但仍開機狀態：} 使用率低但仍開機的伺服器，可作為整併或睡眠狀態建議的候選。
\item \textbf{功率--溫度衝突：} 即使機櫃仍有電氣容量，只要熱餘裕不足，就不應再接收額外工作負載。
\end{itemize}

機櫃利用率不應只由功率定義。即使某個機櫃的電氣容量仍可用，但若其熱餘裕不足，也不應接收更多工作負載。因此，建議使用一個聯合機櫃狀態指標，同時考慮功率容量、作用中伺服器數量、預測最高溫度與熱點風險。

\subsection{基於數位孿生的強化學習}

在數位孿生或模擬器被驗證之後，強化學習可以作為決策最佳化層加入。在本報告中，強化學習不是 XGBoost/LightGBM 預測的替代品，而是使用已預測與重建的熱狀態來選擇預防行動。

強化學習的形式可定義如下：
\begin{itemize}
\item \textbf{狀態：} 預測伺服器溫度、預測機櫃最高溫度、熱點風險分數、機櫃區域功率、風扇速度、PDU 功率、重建熱場或數位孿生潛在狀態，以及工作負載狀態圖。
\item \textbf{動作：} 工作負載遷移、工作負載配置、工作負載降頻、伺服器睡眠/閒置建議、冷卻設定點調整、氣流調整，或風扇速度調整。
\item \textbf{獎勵：} 降低能源使用、降低 PUE 代理指標、降低熱點風險，以及提升機櫃利用率可給予正獎勵；違反熱限制、高溫不均勻性、SLA/效能退化、過度遷移與快速控制震盪則給予負懲罰。
\item \textbf{安全限制：} 最大 CPU/GPU 溫度、最大機櫃溫度、最小熱餘裕、最大遷移率、允許的控制範圍，以及動作平滑限制。
\end{itemize}

強化學習的實作應遵循分階段路徑：
\begin{enumerate}
\item 由遙測資料、機櫃中繼資料與重建熱場建立並校準數位孿生。
\item 驗證數位孿生能否重現工作負載或冷卻變化後的觀察溫度反應。
\item 在數位孿生內使用安全動作範圍離線訓練強化學習代理。
\item 在相同模擬工作負載下，比較規則式、啟發式、數學最佳化與強化學習策略。
\item 讓強化學習代理先以 shadow mode 執行，也就是只提供行動建議，但不控制實體系統。
\item 只有在安全過濾器與回復規則被定義後，才啟用閉迴路控制。
\end{enumerate}

在動作空間設計方面，當動作是離散的，例如選擇要進行工作負載遷移或睡眠建議的伺服器時，DQN 類型方法是合適的~\cite{yuan_2025}。若動作同時包含離散與連續項，例如選擇目標伺服器，同時調整氣流或冷卻設定點，則 DeepEE 中使用的 PADQN 方向較為適合~\cite{ran_2019}。這個分階段強化學習層直接支援預防，因為它能在過熱發生之前評估行動後果。

\section{討論與開放議題}

已回顧的論文共同顯示，沒有單一模型能解決完整問題。XGBoost 與 LightGBM 適合結構化遙測資料，但機櫃層級預測也需要空間背景。GRU/LSTM 模型可以處理時間延遲，但需要可靠的歷史軌跡。U-Net 與隱式神經表示法能夠重建熱場，但需要感測器座標與訓練熱圖。數位孿生是預測與控制之間自然的橋梁，因為它能將重建熱狀態視覺化，並在部署前評估冷卻或工作負載行動的影響。強化學習能最佳化冷卻與工作負載決策，但應在安全模擬器、數位孿生，或具規則限制的環境可用後，作為未來決策層提出。

本專案的主要開放議題如下：
\begin{itemize}
\item \textbf{資料完整性：} CPU 使用率與多個 Node Exporter 指標雖然有用，但在可用的 CortexDC 機櫃文件中並不一致存在。
\item \textbf{感測器座標：} 熱重建需要感測器位置、機櫃幾何與一致的空間網格。
\item \textbf{鄰近定義：} 鄰近機櫃/伺服器溫度應根據實體相鄰關係、氣流路徑或機櫃區域來定義。
\item \textbf{真實標籤：} 機櫃層級標籤，例如機櫃最高溫度或熱點位置，需要足夠的感測器或經驗證的熱圖。
\item \textbf{控制安全性：} 工作負載遷移、風扇控制與冷卻設定點變更，在自動化之前需要操作限制，最好先在模擬器或數位孿生中評估。
\end{itemize}

\subsection{驗證策略}

第一個驗證步驟應先評估遙測資料完整性，而不是直接評估模型準確率。每個欄位都應檢查缺失率、取樣間隔、單位一致性，以及是否能對應到正確的伺服器、插座、機櫃與 U 位置。在可用特徵集合固定之後，XGBoost 與 LightGBM 可使用時間感知的訓練/測試切分進行評估，避免用未來樣本預測過去樣本。建議指標包括平均絕對誤差、均方根誤差、最高溫度誤差、熱點分類精確率/召回率，以及缺失資料穩健性。

機櫃層級模型也應與伺服器層級模型分開驗證。伺服器層級驗證檢查是否能預測個別機器的 CPU、GPU 或進氣溫度。機櫃層級驗證則檢查聚合特徵與空間特徵，是否相較於非空間基準模型，能改善機櫃平均溫度、機櫃最高溫度或熱點風險預測。

\subsection{實作路線圖}

建議的實作順序如下：
\begin{enumerate}
\item 建立遙測稽核工具，報告 Redfish、PDU、環境與機櫃中繼資料欄位在時間上的可用情況。
\item 使用已證實的功率、溫度、風扇、PDU 與機櫃中繼資料特徵，實作保守基準模型。
\item 訓練 XGBoost/LightGBM 模型進行伺服器層級與機櫃層級預測，並與簡單基準模型比較。
\item 加入空間特徵，例如機櫃區域、加權 U 位置，以及鄰近熱/功率特徵。
\item 只有在穩定的歷史遙測資料與空間標籤可用後，才加入時序模型、熱重建與數位孿生視覺化層。
\item 先將預測用於警告與建議；在任何閉迴路工作負載或冷卻控制之前，先於數位孿生或具規則限制的模擬器中評估強化學習決策。
\end{enumerate}

\section{結論}

本調查整理了可支援資料中心熱監測、機櫃層級溫度預測與功率感知利用的 AI/ML 模型。XGBoost 與 LightGBM 適合作為 CortexDC 已驗證結構化遙測資料的基準模型，因為它們能從伺服器功率、進氣溫度、CPU 溫度、風扇速度、PDU 功率與機櫃中繼資料中學習非線性關係。然而，機櫃層級預測必須包含空間特徵，例如機櫃 ID、U 位置、機櫃區域、鄰近伺服器/機櫃溫度，以及機櫃功率分布。

對於未來延伸，GRU/LSTM 模型可以建模熱延遲；U-Net 或隱式神經表示法可以重建稀疏熱場；而數位孿生可以提供視覺化與安全行動評估。強化學習則應被定位為較後期的決策最佳化層，用於聯合工作負載與冷卻控制，而不是第一個預測模型。

所提出的架構刻意保持保守。它只宣稱使用 Redfish、PDU 與機櫃文件已證實的資料欄位，並將 CPU 使用率與 Node Exporter 指標視為選用，除非它們在遙測紀錄中穩定存在。這使得此架構可作為第一個基準模型實作，同時也保留了通往更完整數位孿生與最佳化系統的明確路徑。

\begin{thebibliography}{00}

\bibitem{lin_review_2024}
J. Lin, W. Lin, H. Huang, W. Lin, and K. Li, ``Thermal modeling and thermal-aware energy saving methods for cloud data centers: A review,'' \textit{IEEE Transactions on Sustainable Computing}, vol. 9, no. 3, pp. 571--590, May--Jun. 2024, doi: 10.1109/TSUSC.2023.3346332.

\bibitem{zhang_2018}
K. Zhang, A. Guliani, S. Ogrenci-Memik, G. Memik, K. Yoshii, R. Sankaran, and P. Beckman, ``Machine learning-based temperature prediction for runtime thermal management across system components,'' \textit{IEEE Transactions on Parallel and Distributed Systems}, vol. 29, no. 2, pp. 405--419, Feb. 2018, doi: 10.1109/TPDS.2017.2732951.

\bibitem{sano_2018}
K. Sano, H. Shimizu, Y. Kondo, and T. Fujimoto, ``Improving accuracy of temperature distribution and energy-saving technology of air conditioners in data centers,'' \textit{IEEE Transactions on Components, Packaging and Manufacturing Technology}, vol. 8, no. 5, pp. 811--817, May 2018, doi: 10.1109/TCPMT.2018.2811545.

\bibitem{xwang_sensor_2013}
X. Wang, X. Wang, G. Xing, J. Chen, C.-X. Lin, and Y. Chen, ``Intelligent sensor placement for hot server detection in data centers,'' \textit{IEEE Transactions on Parallel and Distributed Systems}, vol. 24, no. 8, pp. 1577--1588, Aug. 2013, doi: 10.1109/TPDS.2012.254.

\bibitem{wang_self_healing_2024}
S. Wang, Y. Kang, Y. Xu, C. Ma, H. Wang, and W. Wu, ``Data center temperature prediction and management based on a two-stage self-healing model,'' \textit{Simulation Modelling Practice and Theory}, vol. 132, Art. no. 102883, Apr. 2024, doi: 10.1016/j.simpat.2023.102883.

\bibitem{yuan_2025}
H. Yuan, Z. Zhang, D. Yang, T. Xue, D. Wen, and G. Yao, ``An intelligent thermal management strategy for a data center prototype based on digital twin technology,'' \textit{Applied Sciences}, vol. 15, no. 14, Art. no. 7675, Jul. 2025, doi: 10.3390/app15147675.

\bibitem{andersson_2023}
T. R. Andersson, W. P. Bruinsma, S. Markou, J. Requeima, A. Coca-Castro, A. Vaughan, A.-L. Ellis, M. A. Lazzara, D. Jones, J. S. Hosking, and R. E. Turner, ``Environmental sensor placement with convolutional Gaussian neural processes,'' \textit{Environmental Data Science}, vol. 2, Art. no. e32, 2023, doi: 10.1017/eds.2023.22.

\bibitem{nisce_2023}
I. Nisce, X. Jiang, and S. P. Vishnu, ``Machine learning based thermal prediction for energy-efficient cloud computing,'' in \textit{Proc. IEEE 20th Consumer Communications \& Networking Conf. (CCNC)}, Las Vegas, NV, USA, 2023, pp. 624--627, doi: 10.1109/CCNC51644.2023.10060079.

\bibitem{ilager_2021}
S. Ilager, K. Ramamohanarao, and R. Buyya, ``Thermal prediction for efficient energy management of clouds using machine learning,'' \textit{IEEE Transactions on Parallel and Distributed Systems}, vol. 32, no. 5, pp. 1044--1056, May 2021, doi: 10.1109/TPDS.2020.3040800.

\bibitem{lin_ate_2022}
J. Lin, W. Lin, W. Lin, J. Wang, and H. Jiang, ``Thermal prediction for air-cooled data center using data driven-based model,'' \textit{Applied Thermal Engineering}, vol. 217, Art. no. 119207, Nov. 2022, doi: 10.1016/j.applthermaleng.2022.119207.

\bibitem{gill_thermosim_2020}
S. S. Gill, S. Tuli, A. N. Toosi, F. Cuadrado, P. Garraghan, R. Bahsoon, H. Lutfiyya, R. Sakellariou, O. Rana, S. Dustdar, and R. Buyya, ``ThermoSim: Deep learning based framework for modeling and simulation of thermal-aware resource management for cloud computing environments,'' \textit{Journal of Systems and Software}, vol. 166, Art. no. 110596, Aug. 2020, doi: 10.1016/j.jss.2020.110596.

\bibitem{sabathiel_2024}
S. Sabathiel, H. Sanchis-Alepuz, A. S. Wilson, J. Reynvaan, and M. Stipsitz, ``Neural network-based reconstruction of steady-state temperature systems with unknown material composition,'' \textit{Scientific Reports}, vol. 14, Art. no. 22265, Sep. 2024, doi: 10.1038/s41598-024-73380-1.

\bibitem{peng_2022}
X. Peng, X. Li, Z. Gong, X. Zhao, and W. Yao, ``A deep learning method based on partition modeling for reconstructing temperature field,'' \textit{International Journal of Thermal Sciences}, vol. 182, Art. no. 107802, Dec. 2022, doi: 10.1016/j.ijthermalsci.2022.107802.

\bibitem{luo_2024}
X. Luo, W. Xu, B. T. Nadiga, Y. Ren, and S. Yoo, ``Continuous field reconstruction from sparse observations with implicit neural networks.'' in \textit{Proc. 12th International Conference on Learning Representations (ICLR)}, 2024.

\bibitem{zapater_2015}
M. Zapater, O. Tuncer, J. L. Ayala, J. M. Moya, K. Vaidyanathan, K. C. Gross, and A. K. Coskun, ``Leakage-aware cooling management for improving server energy efficiency,'' \textit{IEEE Transactions on Parallel and Distributed Systems}, vol. 26, no. 10, pp. 2764--2777, Oct. 2015, doi: 10.1109/TPDS.2014.2361519.

\bibitem{chan_2019}
C. S. Chan, A. S. Akyurek, B. Aksanli, and T. S. Rosing, ``Optimal performance-aware cooling on enterprise servers,'' \textit{IEEE Transactions on Computer-Aided Design of Integrated Circuits and Systems}, vol. 38, no. 9, pp. 1689--1702, Sep. 2019, doi: 10.1109/TCAD.2018.2855122.

\bibitem{ran_2019}
Y. Ran, H. Hu, X. Zhou, and Y. Wen, ``DeepEE: Joint optimization of job scheduling and cooling control for data center energy efficiency using deep reinforcement learning,'' in \textit{Proc. IEEE 39th International Conference on Distributed Computing Systems (ICDCS)}, Dallas, TX, USA, 2019, pp. 645--655, doi: 10.1109/ICDCS.2019.00070.

\bibitem{cortexdc_redfish_ref}
CortexDC, ``CortexDC server Redfish monitoring data reference,'' internal technical reference, 2026.

\bibitem{cortexdc_pdu_ref}
CortexDC, ``CortexDC PDU monitoring data and communication protocol reference,'' internal technical reference, 2026.

\bibitem{cortexdc_datasheet_2026}
CortexDC Management System, ``CortexDC data center inventory and rack infrastructure documentation,''internal technical report, May 21, 2026.

\end{thebibliography}

\end{document}